\documentclass[11pt,a4paper]{article}

\usepackage[utf8]{inputenc}
\usepackage[T1]{fontenc}
\usepackage{geometry}
\usepackage{hyperref}
\usepackage{enumitem}
\usepackage{fancyhdr}
\usepackage{titlesec}
\usepackage{longtable}
\usepackage{booktabs}
\usepackage{listings}
\usepackage{xcolor}

\geometry{margin=1in}
\hypersetup{colorlinks=true, linkcolor=blue, urlcolor=blue}

\lstset{
  basicstyle=\ttfamily\small,
  breaklines=true,
  frame=single,
  backgroundcolor=\color{gray!10}
}

\pagestyle{fancy}
\fancyhf{}
\fancyhead[L]{\textbf{AICT API \& Schema Specification}}
\fancyhead[R]{\thepage}
\renewcommand{\headrulewidth}{0.4pt}

\title{\textbf{AICT Multi-Agent Management Platform} \\ \Large API \& Schema Specification}
\author{AICT Project}
\date{\today}

\begin{document}
\maketitle
\tableofcontents
\newpage

%% ============================================================
\section{System Architecture Overview}
%% ============================================================

\subsection{High-Level Components}

\begin{enumerate}
  \item \textbf{React Frontend} --- SPA with three views: Chat, Kanban, Spec File Manager
  \item \textbf{Python Backend (FastAPI)} --- REST API + WebSocket server
  \item \textbf{PostgreSQL Database} --- All persistent data, ticket queue, pgvector embeddings
  \item \textbf{E2B Sandboxes} --- Isolated runtime environments for agents
  \item \textbf{Git Repositories} --- Spec folder repo (platform-managed) + Project code repo (user's repo)
  \item \textbf{LaTeX Container} --- Docker container for pdflatex compilation
\end{enumerate}

\subsection{Data Flow Overview}

\begin{enumerate}
  \item \textbf{Client $\to$ Backend}: REST API calls (CRUD), WebSocket for chat messages
  \item \textbf{Backend $\to$ Client}: WebSocket pushes (chat responses, Kanban updates, agent activity)
  \item \textbf{Backend $\to$ Agents}: Ticket dispatch via Agent Tooling API; agent invocation via E2B
  \item \textbf{Agents $\to$ Backend}: Tool calls (read/write files, Git operations, ticket creation, Kanban updates)
  \item \textbf{Agents $\to$ Git}: Clone, branch, commit, push, create PR (project repo)
  \item \textbf{Backend $\to$ LaTeX Container}: Compile .tex files, return .pdf
\end{enumerate}

\subsection{Component Diagram}

\begin{verbatim}
+------------------+       WebSocket / REST       +------------------+
|  React Frontend  | <------------------------->  |  FastAPI Backend |
+------------------+                              +------------------+
                                                         |
                    +------------------------------------+------------------------------------+
                    |                    |                    |                    |
              +-----v-----+        +-----v-----+        +-----v-----+        +-----v-----+
              | PostgreSQL|        |    E2B    |        |    Git    |        |   LaTeX   |
              |  Database |        | Sandboxes |        |   Repos   |        | Container |
              +-----------+        +-----------+        +-----------+        +-----------+
\end{verbatim}

%% ============================================================
\section{Database Schema}
%% ============================================================

All tables use UUID primary keys. Timestamps are \texttt{TIMESTAMPTZ}.

\subsection{Table: projects}

Single project per instance (for initial release).

\begin{lstlisting}
CREATE TABLE projects (
  id              UUID PRIMARY KEY DEFAULT gen_random_uuid(),
  name            VARCHAR(255) NOT NULL,
  description     TEXT,
  spec_repo_path  VARCHAR(512) NOT NULL,  -- local path to spec Git repo
  code_repo_url   VARCHAR(512) NOT NULL,  -- remote URL of project code repo
  code_repo_path  VARCHAR(512) NOT NULL,  -- local clone path
  created_at      TIMESTAMPTZ DEFAULT NOW(),
  updated_at      TIMESTAMPTZ DEFAULT NOW()
);
\end{lstlisting}

\subsection{Table: agents}

Tracks all agent instances.

\begin{lstlisting}
CREATE TABLE agents (
  id              UUID PRIMARY KEY DEFAULT gen_random_uuid(),
  project_id      UUID REFERENCES projects(id) ON DELETE CASCADE,
  role            VARCHAR(50) NOT NULL,  -- 'gm', 'om', 'engineer'
  level           VARCHAR(10),           -- NULL for GM/OM; 'L1' or 'L2' for engineers
  display_name    VARCHAR(100) NOT NULL, -- e.g. 'GM', 'OM-1', 'Engineer-3'
  model           VARCHAR(100) NOT NULL, -- e.g. 'gemini-3-pro', 'claude-4.5-opus'
  status          VARCHAR(20) DEFAULT 'sleeping', -- 'sleeping', 'active', 'busy'
  current_task_id UUID,                  -- FK to tasks, nullable
  sandbox_id      VARCHAR(255),          -- E2B sandbox ID, NULL if not running
  sandbox_persist BOOLEAN DEFAULT FALSE, -- TRUE for GM/OM, FALSE for engineers
  priority        INT DEFAULT 2,         -- 0=GM, 1=OM, 2=Engineer
  created_at      TIMESTAMPTZ DEFAULT NOW(),
  updated_at      TIMESTAMPTZ DEFAULT NOW()
);
\end{lstlisting}

\subsection{Table: tasks}

Kanban cards.

\begin{lstlisting}
CREATE TABLE tasks (
  id              UUID PRIMARY KEY DEFAULT gen_random_uuid(),
  project_id      UUID REFERENCES projects(id) ON DELETE CASCADE,
  title           VARCHAR(255) NOT NULL,
  description     TEXT,
  status          VARCHAR(50) DEFAULT 'backlog',
    -- 'backlog', 'specifying', 'assigned', 'in_progress', 'in_review', 'done'
  critical        INT DEFAULT 8,         -- 0-10, 0=most critical
  urgent          INT DEFAULT 8         -- 0-10, 0=most urgent
  assigned_agent  UUID REFERENCES agents(id) ON DELETE SET NULL,
  module_path     VARCHAR(512),          -- file/folder scope for engineer
  git_branch      VARCHAR(255),          -- branch name if applicable
  pr_url          VARCHAR(512),          -- PR URL once created
  parent_task_id  UUID REFERENCES tasks(id) ON DELETE SET NULL, -- for subtasks
  created_by      UUID REFERENCES agents(id),  -- which agent created this task
  created_at      TIMESTAMPTZ DEFAULT NOW(),
  updated_at      TIMESTAMPTZ DEFAULT NOW()
);
\end{lstlisting}

\subsection{Table: tickets}

Agent-to-agent communication queue.

\begin{lstlisting}
CREATE TABLE tickets (
  id              UUID PRIMARY KEY DEFAULT gen_random_uuid(),
  project_id      UUID REFERENCES projects(id) ON DELETE CASCADE,
  from_agent      UUID REFERENCES agents(id) ON DELETE CASCADE,
  to_agent        UUID REFERENCES agents(id) ON DELETE CASCADE,
  header          VARCHAR(255) NOT NULL,
  ticket_type     VARCHAR(50) NOT NULL,
    -- 'task_assignment', 'question', 'help', 'issue'
  critical        INT DEFAULT 8,
  urgent          INT DEFAULT 8,
  status          VARCHAR(20) DEFAULT 'open',  -- 'open', 'closed'
  created_at      TIMESTAMPTZ DEFAULT NOW(),
  closed_at       TIMESTAMPTZ,
  closed_by       UUID REFERENCES agents(id)
);
\end{lstlisting}

\subsection{Table: ticket\_messages}

Conversation within a ticket.

\begin{lstlisting}
CREATE TABLE ticket_messages (
  id              UUID PRIMARY KEY DEFAULT gen_random_uuid(),
  ticket_id       UUID REFERENCES tickets(id) ON DELETE CASCADE,
  from_agent      UUID REFERENCES agents(id) ON DELETE CASCADE,
  content         TEXT NOT NULL,
  created_at      TIMESTAMPTZ DEFAULT NOW()
);
\end{lstlisting}

\subsection{Table: chat\_messages}

Client $\leftrightarrow$ GM conversation.

\begin{lstlisting}
CREATE TABLE chat_messages (
  id              UUID PRIMARY KEY DEFAULT gen_random_uuid(),
  project_id      UUID REFERENCES projects(id) ON DELETE CASCADE,
  role            VARCHAR(20) NOT NULL,  -- 'user', 'gm'
  content         TEXT NOT NULL,
  attachments     JSONB,                 -- array of file references
  created_at      TIMESTAMPTZ DEFAULT NOW()
);
\end{lstlisting}

\subsection{Table: spec\_files}

Metadata for files in the spec Git repo. Actual content is in Git.

\begin{lstlisting}
CREATE TABLE spec_files (
  id              UUID PRIMARY KEY DEFAULT gen_random_uuid(),
  project_id      UUID REFERENCES projects(id) ON DELETE CASCADE,
  path            VARCHAR(512) NOT NULL, -- relative path in spec repo
  file_type       VARCHAR(20),           -- 'tex', 'pdf', 'image', 'other'
  last_commit     VARCHAR(40),           -- Git commit SHA
  last_modified   TIMESTAMPTZ,
  UNIQUE(project_id, path)
);
\end{lstlisting}

\subsection{Table: embeddings}

pgvector table for RAG.

\begin{lstlisting}
CREATE EXTENSION IF NOT EXISTS vector;

CREATE TABLE embeddings (
  id              UUID PRIMARY KEY DEFAULT gen_random_uuid(),
  project_id      UUID REFERENCES projects(id) ON DELETE CASCADE,
  source_type     VARCHAR(50) NOT NULL,  -- 'spec_file', 'chat', 'code'
  source_id       UUID,                  -- FK to source table (nullable for flexibility)
  source_path     VARCHAR(512),          -- file path or identifier
  chunk_index     INT DEFAULT 0,         -- chunk number within source
  content         TEXT NOT NULL,         -- the text chunk
  embedding       vector(1536),          -- OpenAI ada-002 dimension; adjust as needed
  created_at      TIMESTAMPTZ DEFAULT NOW()
);

CREATE INDEX ON embeddings USING ivfflat (embedding vector_cosine_ops) WITH (lists = 100);
\end{lstlisting}

\subsection{Table: agent\_activity\_log}

For the live activity feed.

\begin{lstlisting}
CREATE TABLE agent_activity_log (
  id              UUID PRIMARY KEY DEFAULT gen_random_uuid(),
  project_id      UUID REFERENCES projects(id) ON DELETE CASCADE,
  agent_id        UUID REFERENCES agents(id) ON DELETE CASCADE,
  activity        VARCHAR(255) NOT NULL, -- e.g. 'writing tests', 'reviewing PR'
  task_id         UUID REFERENCES tasks(id) ON DELETE SET NULL,
  created_at      TIMESTAMPTZ DEFAULT NOW()
);
\end{lstlisting}

%% ============================================================
\section{REST API Endpoints}
%% ============================================================

Base URL: \texttt{/api/v1}

\subsection{Project}

\begin{longtable}{p{2.5cm} p{4cm} p{7cm}}
\toprule
\textbf{Method} & \textbf{Endpoint} & \textbf{Description} \\
\midrule
\endhead
GET    & /project              & Get current project info \\
PUT    & /project              & Update project settings \\
POST   & /project/init         & Initialize new project (spec repo + clone code repo) \\
\bottomrule
\end{longtable}

\subsection{Chat (Client $\leftrightarrow$ GM)}

\begin{longtable}{p{2.5cm} p{4cm} p{7cm}}
\toprule
\textbf{Method} & \textbf{Endpoint} & \textbf{Description} \\
\midrule
\endhead
GET    & /chat/messages        & Get chat history (paginated) \\
POST   & /chat/messages        & Send message to GM (triggers GM agent) \\
GET    & /chat/status          & Get GM status (available/busy) \\
\bottomrule
\end{longtable}

\subsection{Kanban / Tasks}

\begin{longtable}{p{2.5cm} p{4cm} p{7cm}}
\toprule
\textbf{Method} & \textbf{Endpoint} & \textbf{Description} \\
\midrule
\endhead
GET    & /tasks                & List all tasks (filterable by status, agent, module) \\
GET    & /tasks/:id            & Get single task details \\
POST   & /tasks                & Create task (user or GM) \\
PUT    & /tasks/:id            & Update task (status, assignment, etc.) \\
DELETE & /tasks/:id            & Delete task \\
GET    & /tasks/:id/subtasks   & Get subtasks of a task \\
\bottomrule
\end{longtable}

\subsection{Agents}

\begin{longtable}{p{2.5cm} p{4cm} p{7cm}}
\toprule
\textbf{Method} & \textbf{Endpoint} & \textbf{Description} \\
\midrule
\endhead
GET    & /agents               & List all agents \\
GET    & /agents/:id           & Get agent details \\
GET    & /agents/:id/activity  & Get recent activity for agent \\
GET    & /agents/activity      & Get live activity feed (all agents) \\
\bottomrule
\end{longtable}

\subsection{Tickets}

\begin{longtable}{p{2.5cm} p{4cm} p{7cm}}
\toprule
\textbf{Method} & \textbf{Endpoint} & \textbf{Description} \\
\midrule
\endhead
GET    & /tickets              & List tickets (filterable) \\
GET    & /tickets/:id          & Get ticket with messages \\
\bottomrule
\end{longtable}

Note: Tickets are created/managed by agents via the Agent Tooling API, not by the client directly.

\subsection{Spec Files}

\begin{longtable}{p{2.5cm} p{4cm} p{7cm}}
\toprule
\textbf{Method} & \textbf{Endpoint} & \textbf{Description} \\
\midrule
\endhead
GET    & /specs                & List spec folder contents (tree) \\
GET    & /specs/*path          & Get file content (text) or download (binary) \\
PUT    & /specs/*path          & Create or update file \\
DELETE & /specs/*path          & Delete file or folder \\
POST   & /specs/folder         & Create folder \\
POST   & /specs/compile        & Compile .tex file, return PDF \\
GET    & /specs/history/*path  & Get version history for file \\
GET    & /specs/diff           & Get diff between commits \\
\bottomrule
\end{longtable}

%% ============================================================
\section{WebSocket API}
%% ============================================================

Endpoint: \texttt{/ws}

All messages are JSON with a \texttt{type} field.

\subsection{Client $\to$ Server}

\begin{lstlisting}
// Send chat message to GM
{
  "type": "chat_message",
  "content": "I want to add a login feature",
  "attachments": []  // optional file references
}

// Subscribe to updates (on connect)
{
  "type": "subscribe",
  "channels": ["chat", "kanban", "activity"]
}
\end{lstlisting}

\subsection{Server $\to$ Client}

\begin{lstlisting}
// GM chat response (streamed)
{
  "type": "chat_message",
  "role": "gm",
  "content": "Sure, let me clarify...",
  "message_id": "uuid",
  "is_complete": false  // true when stream ends
}

// GM status change
{
  "type": "gm_status",
  "status": "busy",
  "busy_with": "talking to OM-1"
}

// Kanban task update
{
  "type": "task_update",
  "task": { ...task object... }
}

// Task created
{
  "type": "task_created",
  "task": { ...task object... }
}

// Task deleted
{
  "type": "task_deleted",
  "task_id": "uuid"
}

// Agent activity
{
  "type": "agent_activity",
  "agent_id": "uuid",
  "agent_name": "Engineer-2",
  "activity": "writing unit tests",
  "task_id": "uuid"
}

// Agent status change
{
  "type": "agent_status",
  "agent_id": "uuid",
  "status": "active"  // or 'sleeping', 'busy'
}

// Spec file change (e.g., GM or client edited)
{
  "type": "spec_file_changed",
  "path": "specs/backend/auth.tex",
  "action": "modified"  // or 'created', 'deleted'
}
\end{lstlisting}

%% ============================================================
\section{Agent Tooling API}
%% ============================================================

These are the tools available to agents inside their E2B sandbox. Agents call these via HTTP to the backend (internal endpoint, not exposed to client).

Internal base URL: \texttt{http://backend:8000/internal/agent}

All requests include \texttt{X-Agent-ID} header.

\subsection{Ticket Operations}

\begin{longtable}{p{2.5cm} p{5cm} p{6cm}}
\toprule
\textbf{Method} & \textbf{Endpoint} & \textbf{Description} \\
\midrule
\endhead
POST   & /tickets                   & Create ticket to another agent \\
GET    & /tickets/inbox             & Get pending tickets for this agent \\
POST   & /tickets/:id/reply         & Reply to a ticket \\
POST   & /tickets/:id/close         & Close a ticket (higher priority only) \\
\bottomrule
\end{longtable}

\textbf{Create ticket request body:}
\begin{lstlisting}
{
  "to_agent_role": "om",       // or specific agent_id
  "header": "Need clarification on auth module",
  "ticket_type": "question",   // 'task_assignment', 'question', 'help', 'issue'
  "critical": 5,
  "urgent": 3,
  "initial_message": "The spec says OAuth but doesn't specify provider..."
}
\end{lstlisting}

\subsection{Task / Kanban Operations}

\begin{longtable}{p{2.5cm} p{5cm} p{6cm}}
\toprule
\textbf{Method} & \textbf{Endpoint} & \textbf{Description} \\
\midrule
\endhead
GET    & /tasks                     & List tasks (filtered by assignment, status) \\
GET    & /tasks/:id                 & Get task details \\
POST   & /tasks                     & Create task (OM/GM only) \\
PUT    & /tasks/:id                 & Update task (status, description, etc.) \\
PUT    & /tasks/:id/assign          & Assign task to agent (OM only) \\
POST   & /tasks/:id/subtasks        & Create subtask \\
\bottomrule
\end{longtable}

\subsection{Git Operations}

Wrapped through GitService with guardrails.

\begin{longtable}{p{2.5cm} p{5cm} p{6cm}}
\toprule
\textbf{Method} & \textbf{Endpoint} & \textbf{Description} \\
\midrule
\endhead
POST   & /git/branch                & Create new branch \\
POST   & /git/checkout              & Checkout branch \\
POST   & /git/commit                & Commit changes \\
POST   & /git/push                  & Push to remote \\
POST   & /git/pull                  & Pull from remote \\
POST   & /git/pr                    & Create pull request \\
POST   & /git/merge                 & Merge PR (OM only) \\
GET    & /git/status                & Get git status \\
GET    & /git/diff                  & Get diff \\
GET    & /git/log                   & Get commit log \\
\bottomrule
\end{longtable}

\textbf{Blocked commands} (return 403):
\begin{itemize}
  \item \texttt{git rebase}
  \item \texttt{git reset --hard}
  \item \texttt{git push --force}
  \item \texttt{git commit --amend} (on pushed commits)
  \item \texttt{git filter-branch}
  \item \texttt{git reflog expire}
\end{itemize}

\subsection{File Operations}

\begin{longtable}{p{2.5cm} p{5cm} p{6cm}}
\toprule
\textbf{Method} & \textbf{Endpoint} & \textbf{Description} \\
\midrule
\endhead
GET    & /files/read                & Read file content \\
POST   & /files/write               & Write file (scope-checked) \\
DELETE & /files/delete              & Delete file (scope-checked) \\
GET    & /files/list                & List directory contents \\
POST   & /files/search              & Search files by pattern \\
\bottomrule
\end{longtable}

\textbf{Scope enforcement:}
\begin{itemize}
  \item Engineers: can only write to their assigned \texttt{module\_path}
  \item OM: can write to spec files and task definitions
  \item GM: can write to spec folder and project config
\end{itemize}

\subsection{Spec Folder Operations (GM/OM)}

\begin{longtable}{p{2.5cm} p{5cm} p{6cm}}
\toprule
\textbf{Method} & \textbf{Endpoint} & \textbf{Description} \\
\midrule
\endhead
GET    & /specs/read/*path          & Read spec file \\
POST   & /specs/write/*path         & Write spec file (auto-commits) \\
DELETE & /specs/delete/*path        & Delete spec file (auto-commits) \\
POST   & /specs/compile/*path       & Compile .tex to PDF \\
\bottomrule
\end{longtable}

\subsection{RAG / Knowledge Retrieval}

\begin{longtable}{p{2.5cm} p{5cm} p{6cm}}
\toprule
\textbf{Method} & \textbf{Endpoint} & \textbf{Description} \\
\midrule
\endhead
POST   & /rag/query                 & Semantic search over project knowledge \\
POST   & /rag/index                 & Index new content (triggered automatically) \\
\bottomrule
\end{longtable}

\textbf{Query request:}
\begin{lstlisting}
{
  "query": "How does the auth module handle sessions?",
  "source_types": ["spec_file", "code"],  // optional filter
  "top_k": 5
}
\end{lstlisting}

\subsection{Agent Lifecycle}

\begin{longtable}{p{2.5cm} p{5cm} p{6cm}}
\toprule
\textbf{Method} & \textbf{Endpoint} & \textbf{Description} \\
\midrule
\endhead
POST   & /agent/status              & Update own status (active/busy/sleeping) \\
POST   & /agent/activity            & Log activity (for live feed) \\
POST   & /agent/sleep               & Declare intent to sleep (close sandbox if engineer) \\
\bottomrule
\end{longtable}

%% ============================================================
\section{E2B Sandbox Lifecycle}
%% ============================================================

\subsection{GM and Operation Master}

\begin{enumerate}
  \item Sandbox created on project initialization (GM) or when first needed (OM).
  \item Sandbox persists across invocations.
  \item Agent is invoked via ticket or client message.
  \item Agent processes, calls tools, updates status.
  \item Agent declares ``sleeping'' when idle; sandbox remains alive.
  \item Sandbox only destroyed on project close or explicit teardown.
\end{enumerate}

\subsection{Engineer}

\begin{enumerate}
  \item Sandbox created when task is assigned to engineer.
  \item Engineer clones project repo, checks out assigned branch.
  \item Engineer implements, commits, pushes, creates PR.
  \item Engineer updates task status to \texttt{in\_review}.
  \item OM reviews PR:
    \begin{itemize}
      \item If approved: OM merges PR, task status $\to$ \texttt{done}.
      \item If rejected: OM sends ticket with feedback, engineer continues.
    \end{itemize}
  \item When PR merged AND engineer has no more assigned tasks:
    \begin{itemize}
      \item Backend checks \texttt{tasks} table for this agent.
      \item If none pending/in\_progress, sandbox is closed.
    \end{itemize}
  \item If new task assigned later, new sandbox is created.
\end{enumerate}

%% ============================================================
\section{Authentication}
%% ============================================================

Single user per instance. Simple token-based auth.

\begin{lstlisting}
# On startup, backend generates or loads API token from env
API_TOKEN=<random-256-bit-hex>

# Client includes token in all requests
Authorization: Bearer <API_TOKEN>
\end{lstlisting}

WebSocket connection also requires token as query param or first message.

%% ============================================================
\section{LaTeX Compilation}
%% ============================================================

\subsection{Endpoint}

\texttt{POST /api/v1/specs/compile}

\textbf{Request:}
\begin{lstlisting}
{
  "path": "specs/ProductSpec.tex"
}
\end{lstlisting}

\textbf{Response:}
\begin{lstlisting}
{
  "success": true,
  "pdf_path": "specs/ProductSpec.pdf",
  "log": "..."  // pdflatex output, truncated
}
\end{lstlisting}

\subsection{Implementation}

\begin{enumerate}
  \item Backend receives compile request.
  \item Backend runs Docker container: \texttt{docker run --rm -v /spec-repo:/work texlive/texlive pdflatex ...}
  \item Container compiles .tex, outputs .pdf to same directory.
  \item Backend commits the .pdf to spec Git repo.
  \item Backend returns result.
\end{enumerate}

\end{document}
